\documentclass[11pt]{article}
\usepackage[margin=1in]{geometry}
\usepackage{hyperref}
\usepackage{enumitem}
\usepackage{parskip}

\pagestyle{empty}

% Uniform answer block: same indent and spacing under every question.
% Indents the whole block, so wrapped lines stay aligned with the first.
% Leave the argument empty for questions that are not yet answered.
\newcommand{\answer}[1]{%
  \par\vspace{0.35em}%
  \begin{itemize}[label={},leftmargin=1.5em,topsep=0pt,partopsep=0pt,parsep=0pt,itemsep=0pt]
    \item #1
  \end{itemize}%
  \vspace{0.35em}%
}

\begin{document}

\noindent
Candidate's Name: Brendan Premo\\
Phone Number: (518) 588-0855\\
Email: bdanp1998@gmail.com

\vspace{1em}

\begin{enumerate}[leftmargin=*]
    \item Are you currently employed?
    \answer{No, but I am working freelance on some small projects that I am being compensated for by the clients in question.}

    \item How do you ensure data visualization accuracy and integrity?
    \answer{In general, one must be cognizant of the source one is getting the data from and be perceptive enough during an 
    exploratory data analysis (EDA) to find any outliers before making the visualization in the first place. They would 
    also need to double- and triple-check that the calculations were correct (as I've seen a few typos in some of the work
    I had to look over at Soria Analytics) and finally personally consider whether the display is proper for the data, as well
    as get an outsider's opinion on the display to rectify any lack of clarity. For me, if I look at a visualization and
    it doesn't \textit{feel} right, I'll give this list another once-over to see if I missed anything. If I still can't 
    remedy my unease, I'll ask someone to check it over, just to make sure I'm not being overly critical of the work.}

    \item What data visualization tools have you used (e.g., Tableau, Power BI, D3.js)?
    \answer{I primarily used Python and its packages in my schooling (pyplot, seaborn, etc.) and Sigma Computing during my time at 
    Soria Analytics. I've also used many recreationally, such as Tableau, Power BI, Databricks, and, to a lesser extent, Excel for 
    data visualization.}

    \item Can you describe the differences between parametric and non-parametric tests?
    \answer{Parametric and non-parametric tests differ based on the existence of an assumable distribution, for example, a
    normal distribution given by the Central Limit Theorem (CLT) when a sample has size $n\geq30$. Parametric tests rely on a 
    distribution, most commonly a normal distribution, for their tests to work, whereas non-parametric tests don't have the 
    luxury of an underlying distribution and therefore require more nuanced approaches. Metrics are also considered differently; 
    for example, we care a lot about means and standard deviations in parametric tests, but care more about medians and 
    ranks for non-parametric tests. In general, parametric tests have greater statistical power than non-parametric tests, but require 
    assumptions and thus cannot be used on all datasets. Additionally, parametric tests require just the raw data, whereas
    non-parametric tests look at the \textit{ranks} of the data for a proper analysis.}

    \item How do you communicate data findings and recommendations to stakeholders?
    \answer{I think visualization should be doing most of the work; if a chart needs a paragraph of
    explanation before it makes any sense, that's usually a sign the chart needs another pass. There's a school of thought
    that the visual should be the entire deliverable, since anything you write on top of it risks pushing your own
    bias onto whoever is looking at it, and I think that's a fair concern worth designing around. However, I don't think
    the visual can always carry all the understanding an analyst wants to unload. Living in the data day to day means we
    turn up things that aren't obvious at a glance, and there's only ever time to build so many charts, so leaving the whole
    story to whichever visuals happened to get made has its own kind of bias baked into it. Our job is to hand a client 
    actionable insights, and sometimes no amount of polish on a chart is going to get an unseasoned reader to what our
    trained eyes are seeing. So in practice I lead with the visualization and keep a short synopsis next to it: what the
    data shows, what it doesn't, and what I'd want to look at next. I try to keep the wording to what the data actually
    supports and flag where I'm inferring rather than measuring, so a stakeholder can follow the reasoning and push back on
    it instead of just taking my word for it. That saves them time without taking the decision out of their hands.}

    \item Can you explain your experience with machine learning or predictive modeling?
    \answer{During my master's Data Science program, on top of learning and proving the theory mathematically, we were pushed from 
    write the algorithm by hand rather than using prebuilt packages such as scikit learn. This process, while taxing, gave me a 
    deeper understanding of the models I use: both how they work and what use cases they have. After I had finished assignments, 
    I would do them again using the analysis packages to enrich myself further, focusing more on the analysis than the algorithm. 
    Many of my assignments in the Data Science program involved making some kind of predictive model, from EDA to 
    cross-validation to tuning, and a large number of these assignments were in the machine learning curriculum, so I'd say
    I have a lot of experience in my educational background. At Soria Analytics, I made a model to predict the release of new 
    data using Seasonal Autoregressive Integrated Moving Average (SARIMA) that we used for our calendar and newsletter, notifying
    our customers of upcoming data releases.}

    \item How do you handle data confidentiality and security?
    \answer{I try to work inside whatever environment the client already has set up for their data instead of pulling it
    somewhere more convenient for me. At Soria Analytics most of what I touched was government sourced public health data, though some
    of it came to us under a confidentiality agreement, and in those cases the rule of thumb was that the data lives where
    the client says it lives; I'm not making my own copies on my laptop just to make my life easier. If it does have to
    move, it moves encrypted, both on the way there and once it's sitting somewhere. The rest of it is really just habit: I
    only pull the columns I actually need, I don't talk about what I'm finding with anyone outside the client and whoever
    else is on the project with me, and if I'm not sure whether I'm allowed to send something, I ask first rather than
    after. At the end of the day the client is the one on the hook if their data gets out, so their rules beat whatever I
    would have preferred to do.}

    \item How would you approach a complex data analysis problem?
    \answer{I try to break apart big problems into bite-sized pieces, usually by writing up a flow chart or outline on a 
    whiteboard. If I'm working in Python, I like to use Jupyter notebooks and work one cell at a time to check the analysis as I go.
    The notebooks incidentally serve as a paper trail of my work for me or others to later validate the script, and I personally
    feel more comfortable in an environment where I can easily see what's going on with the data.}

    \item How do you handle data security and access control?
    \answer{Generally, people should have the least access they can get away with while still being able to do their job, 
    and that access should come from what role someone has rather than being handed out one person at a time, so when they switch 
    projects or leave you can take it back cleanly instead of hunting around for everything they were given. In practice that 
    means an analyst gets the tables they need and not a blanket admin login, nobody is sharing credentials, and there's a log of who 
    looked at what. At Soria Analytics, this mattered most for our own tooling, i.e. who could edit a pipeline or a dashboard and who 
    could only look at it. With patient data the same idea holds with a lot more on the line, so role-based permissions, two-factor on
    anything touching protected health information, and someone actually going back through the list every so often to see who still 
    needs access, since permissions have a way of piling up if nobody is checking.}

    \item How do you optimize SQL queries for performance?
    \answer{Breaking apart large chunks with nested \texttt{SELECT} subqueries is a good place to start. Being as precise as you
    can with column names, i.e., not using \texttt{SELECT *} but instead \texttt{SELECT column\_name1, column\_name2 ...}, and
    ranges of data, i.e., using \texttt{LIMIT} when testing and using conditional \texttt{WHERE} statements instead of deterministic ones
    (\texttt{WHERE date >= 2020-01-01 AND date < 2021-01-01} vs. \texttt{WHERE date = 2020}) are another two places you can save on
    processing time.}

    \item What certifications or training have you received in data analysis?
    \answer{I have a master's degree in Data Science from SUNY Albany and I was trained in Sigma Computing on the job at 
    Soria Analytics. Most of my experience in other programs is recreational, though I am in the process of getting
    certifications online through LinkedIn Learning.}

    \item What statistical methods have you used in your analysis (e.g., regression, hypothesis testing)?
    \answer{During my schooling, I used everything under the sun: ridge regression, principal component analysis, 
    stochastic gradient descent, Hidden Markov Models and neural networks, to name a few. I have taken courses covering 
    Bayesian statistics, non-parametric methods and other advanced statistical tests, including hypothesis testing
    with various distributions.}

    \item Will you be able to dedicate 20-40 hours a week to the company?
    \answer{Yes. I'd happily work part-time but would prefer a full-time position.}

    \item How long would you expect to work for us if hired and how much will you request per hour if you are hired?
    \answer{My goal after leaving Soria Analytics was to find a stable work environment to fully establish my career
    as a data analyst, and I have a vast love for healthcare data, so I can see myself working here for years. My
    salary expectations are currently within the ballpark of \$75,000, so around \$36 an hour at full-time.}

    \item Can you walk me through a recent data analysis project you worked on?
    \answer{My pride and joy at Soria Analytics was an analysis I did on the Centers for Medicare \& Medicaid Services
    (CMS) Star Ratings of Medicare contracts. The goal was to take the raw data files CMS uploads outlining every contract's 
    rating for each measure of their plan and construct the final star rating for each contract from scratch. Since 
    each measure has varying levels of importance, each measure has varying sway over the contract's overall rating. There 
    were also bonuses given to a contract's rating based on the mean and variance of the measure scores, for example a plan 
    with a high average measure score and a low variance earned a 0.4 boost to the contract's star rating, a middling average 
    and middling variance gave a bonus of 0.1, etc. The contracts were also boosted or penalized by an adjustment called the Categorical Adjustment 
    Index (CAI), which is based on a contract's share of members who are low-income subsidy / dual-eligible or disabled. A 
    contract with a high population of enrollees who meet the above conditions sees their rating increase, whereas a low
    population sees their rating decrease. From these metrics, I needed to calculate the individual measure rating from the 
    raw data files into its proper star using the specialized cut points for each measure, then get each
    contract's raw star rating, add the bonuses, and finally round to the nearest half star for each contract.
    CMS doesn't release all of its methods of calculating the star ratings of contracts, so it isn't feasible to get
    full accuracy on every contract's star rating, but my accuracy was over 92\%, which the client who asked for this
    analysis stated was better than his internal team could get.
    
    The visualizations I used for this analysis were two different Waterfall charts: one for a parent company and one for an 
    individual contract. The contract waterfalls showed the year-to-year change between star ratings of that contract by 
    displaying the major movers of the rating so one could see what made that contract's rating go up or down. The 
    company waterfalls showed which \textit{contracts} caused the most movement in the company's enrollment-weighted
    star rating. This was the real motivation for this analysis, as the files for contract measure data are 
    released prior to CMS's final determination on every contract's star rating, so one can theoretically figure out 
    a parent company's enrollment-weighted average star rating before the official CMS release. To give an example,
    when Humana's H5216 contract, which was roughly 45\% of its Medicare Advantage membership, fell from 4.5 to 3.5 
    stars for 2025, the share of members in 4+ star plans went from 94\% to 25\%, and the stock dropped as much as 22\% in 
    a day, which was a \$4 billion loss of market value. This analysis allows investors to make informed decisions on
    which companies will have a good or bad day at the stock exchange, so it was both rewarding to get an accurate
    result and to create something readily useful for investors.}
\end{enumerate}

\end{document}
